Para $z\ll R$, expande-se $1/(R+z) = (1/R)\cdot 1/(1+z/R)$ usando a série binomial:
\[
\frac{1}{R+z} = \frac{1}{R}\left(1 - \frac{z}{R} + \frac{z^{2}}{R^{2}} - \cdots\right).
\]
Truncando na primeira ordem:
\[
V_g(z) \approx -\frac{GM_E m}{R}\left(1 - \frac{z}{R}\right) = -\frac{GM_E m}{R} + \frac{GM_E m}{R^{2}}\,z.
\]
Como $g_0 = GM_E/R^{2}$ é a aceleração gravitacional na superfície, o segundo termo é exatamente a expressão clássica da energia potencial gravitacional:
\[
V_g(z) \approx -\frac{GM_E m}{R} + m\,g_0\,z = V_g(0) + m\,g_0\,z.\quad\checkmark
\]
Assim, a aproximação linear em $z$ recupera o resultado elementar $\Delta V_g = m\,g_0\,z$, válido para $z \ll R \approx 6400\,\mathrm{km}$.